\begin{solapartat}
\figura[0.3]{sol_fig1.png}
\punts{0,2} + \punts{0,2}

Els dos electrons segueixen una trajectòria circular, ja que la força que hi actua és perpendicular a la seva velocitat, \punts{0,2} els dos electrons gira'n en sentit horari \punts{0,1}
\[ q\ v\ B = m\,\frac{v^2}{R} \ \text{\punts{0,2}} \]
Els electrons descriuen circumferències del mateix radi, ja que les forces tenen el mateix mòdul i els dos electrons tenen la mateixa massa i porten la mateixa velocitat. \punts{0,1}
\end{solapartat}

\begin{solapartat}
\figura[0.35]{sol_fig2.png}
\punts{0,4}

(el camp magnètic també pot anar en sentit contrari). $\vec{B}$ ha de ser paral·lel a la velocitat dels electrons \punts{0,4}, ja que la força serà: $|\vec{F}| = qvB\sin(\phi)$ com que $\phi = 0 \to \vec{F} = 0$ \punts{0,2}
\end{solapartat}
